% GENERATED FILE. Edit canonical lessons, then run npm run generate:books.
% canonical-source-hash: fnv1a64:a4d79593eb6289fa
% canonical-lessons: FA-C447-nabsh, FA-C447-pelak, FA-C447-kodposti, FA-C447-asli, FA-C447-fari

\chapter{Corners and house numbers}
\label{ch:fa-corners-and-house-numbers}
\hlchaptermodality{447}

\begin{chapteropening}
\textbf{By the end of this chapter:} \emph{I can find a building: the street corner, the house number, the postcode, a main road and a side street.}

Complete the last lesson of chapter 447: I can find a building: the street corner, the house number, the postcode, a main road and a side street.
\end{chapteropening}

\section[\texorpdfstring{\fa{نبش} (nabsh)}{نبش (nabsh)}]{\fa{نبش} (nabsh) — a (street) corner}

\label{lesson:FA-C447-nabsh}



\begin{quote}
Before the new one: say the Persian for a distance, then the Persian for to move away.
\end{quote}

\subsection*{You'll want to know: \fa{نبش}}
\textbf{\fa{نبش}} (\emph{nabsh}) — \textquotedblleft{}a (street) corner\textquotedblright{}.

\begin{quote}
\fa{نانوایی} \fa{در} \fa{نبش} \fa{خیابان} \fa{است}.
\end{quote}

\begin{quote}
\emph{Nânvâyi dar nabsh-e khiyâbân ast.}
\end{quote}

\begin{quote}
The bakery is on the corner of the street.
\end{quote}

Used in addresses: \emph{nabsh-e kuche-ye dovvom} — on the corner of the second lane.

\begin{cousinweb}
The origin of \emph{nabsh} is not settled, and EtymDB does not list it. It is the corner where a lane meets a street, and it appears in many Iranian addresses.
\end{cousinweb}

\begin{grammarlens}[title={What you've built}]
One more word for telling, explaining and describing at B1.
\end{grammarlens}

\subsection*{Your turn}
\begin{itemize}
\raggedright
  \item \emph{Say it:} \emph{nabsh}
  \item \emph{Say it:} \emph{nabsh}, once more
  \item \emph{Say it:} a sentence of your own with \emph{nabsh}
  \item \emph{Recall:} say the Persian for a distance, then the Persian for to move away, then say \emph{nabsh} again
\end{itemize}

\begin{tcolorbox}[breakable,colback=teal!4,colframe=teal!35!black,title={Before you move on}]
What does \fa{نبش} mean? (\textquotedblleft{}A (street) corner\textquotedblright{}.) Say it once more.
\end{tcolorbox}
\section[\texorpdfstring{\fa{پلاک} (pelâk)}{پلاک (pelâk)}]{\fa{پلاک} (pelâk) — a house number; a number plate}

\label{lesson:FA-C447-pelak}



\begin{quote}
Before the new one: say the Persian for to move away, then the Persian for a.
\end{quote}

\subsection*{You'll want to know: \fa{پلاک}}
\textbf{\fa{پلاک}} (\emph{pelâk}) — \textquotedblleft{}a house number; a number plate\textquotedblright{}.

\begin{quote}
\fa{خانه} \fa{ما} \fa{پلاک} \fa{دوازده} \fa{است}.
\end{quote}

\begin{quote}
\emph{Khâne-ye mâ pelâk-e davâzdah ast.}
\end{quote}

\begin{quote}
Our house is number twelve.
\end{quote}

\emph{Pelâk-e mâshin} — a car's number plate.

\begin{cousinweb}
\emph{pelâk} was borrowed from French \emph{plaque} 'a plate, a sheet of metal' (EtymDB): the small metal plate with a house number on it.
\end{cousinweb}

\begin{grammarlens}[title={What you've built}]
One more word for telling, explaining and describing at B1.
\end{grammarlens}

\subsection*{Your turn}
\begin{itemize}
\raggedright
  \item \emph{Say it:} \emph{pelâk}
  \item \emph{Say it:} \emph{pelâk}, once more
  \item \emph{Say it:} a sentence of your own with \emph{pelâk}
  \item \emph{Recall:} say the Persian for to move away, then the Persian for a, then say \emph{pelâk} again
\end{itemize}

\begin{tcolorbox}[breakable,colback=teal!4,colframe=teal!35!black,title={Before you move on}]
What does \fa{پلاک} mean? (\textquotedblleft{}A house number; a number plate\textquotedblright{}.) Say it once more.
\end{tcolorbox}
\section[\texorpdfstring{\fa{کدپستی} (kodposti)}{کدپستی (kodposti)}]{\fa{کدپستی} (kodposti) — a postcode (US: a zip code)}

\label{lesson:FA-C447-kodposti}



\begin{quote}
Before the new one: say the Persian for a, then the Persian for a house number.
\end{quote}

\subsection*{You'll want to know: \fa{کدپستی}}
\textbf{\fa{کدپستی}} (\emph{kodposti}) — \textquotedblleft{}a postcode (US: a zip code)\textquotedblright{}.

\begin{quote}
\fa{کدپستی} \fa{را} \fa{روی} \fa{پاکت} \fa{نوشتم}.
\end{quote}

\begin{quote}
\emph{Kodposti râ ru-ye pâkat neveshtam.}
\end{quote}

\begin{quote}
I wrote the postcode on the envelope.
\end{quote}

Iranian postcodes have ten digits.

\begin{cousinweb}
\emph{kodposti} is built from two French loans: \emph{kod} 'code', usually traced to French \emph{code}, and \emph{posti} 'postal', from \emph{post}, usually traced to French \emph{poste}.
\end{cousinweb}

\begin{grammarlens}[title={What you've built}]
One more word for telling, explaining and describing at B1.
\end{grammarlens}

\subsection*{Your turn}
\begin{itemize}
\raggedright
  \item \emph{Say it:} \emph{kodposti}
  \item \emph{Say it:} \emph{kodposti}, once more
  \item \emph{Say it:} a sentence of your own with \emph{kodposti}
  \item \emph{Recall:} say the Persian for a, then the Persian for a house number, then say \emph{kodposti} again
\end{itemize}

\begin{tcolorbox}[breakable,colback=teal!4,colframe=teal!35!black,title={Before you move on}]
What does \fa{کدپستی} mean? (\textquotedblleft{}A postcode (US: a zip code)\textquotedblright{}.) Say it once more.
\end{tcolorbox}
\section[\texorpdfstring{\fa{اصلی} (asli)}{اصلی (asli)}]{\fa{اصلی} (asli) — main, principal}

\label{lesson:FA-C447-asli}



\begin{quote}
Before the new one: say the Persian for a house number, then the Persian for a postcode.
\end{quote}

\subsection*{You'll want to know: \fa{اصلی}}
\textbf{\fa{اصلی}} (\emph{asli}) — \textquotedblleft{}main, principal\textquotedblright{}.

\begin{quote}
\fa{این} \fa{خیابان} \fa{اصلی} \fa{شهر} \fa{است}.
\end{quote}

\begin{quote}
\emph{In khiyâbân-e asli-ye shahr ast.}
\end{quote}

\begin{quote}
This is the main street of the town.
\end{quote}

\emph{Khiyâbân-e asli} — a main road.

\begin{cousinweb}
\emph{asli} was borrowed from Arabic \emph{aṣlī} (EtymDB), an adjective built on \emph{aṣl} 'origin, root'.
\end{cousinweb}

\begin{grammarlens}[title={What you've built}]
One more word for telling, explaining and describing at B1.
\end{grammarlens}

\subsection*{Your turn}
\begin{itemize}
\raggedright
  \item \emph{Say it:} \emph{asli}
  \item \emph{Say it:} \emph{asli}, once more
  \item \emph{Say it:} a sentence of your own with \emph{asli}
  \item \emph{Recall:} say the Persian for a house number, then the Persian for a postcode, then say \emph{asli} again
\end{itemize}

\begin{tcolorbox}[breakable,colback=teal!4,colframe=teal!35!black,title={Before you move on}]
What does \fa{اصلی} mean? (\textquotedblleft{}Main, principal\textquotedblright{}.) Say it once more.
\end{tcolorbox}
\section[\texorpdfstring{\fa{فرعی} (far'i)}{فرعی (far'i)}]{\fa{فرعی} (far'i) — side, minor (of a road); secondary}

\label{lesson:FA-C447-fari}



\begin{quote}
Before the new one: say the Persian for a postcode, then the Persian for main.
\end{quote}

\subsection*{You'll want to know: \fa{فرعی}}
\textbf{\fa{فرعی}} (\emph{far'i}) — \textquotedblleft{}side, minor (of a road); secondary\textquotedblright{}.

\begin{quote}
\fa{از} \fa{خیابان} \fa{اصلی} \fa{وارد} \fa{یک} \fa{خیابان} \fa{فرعی} \fa{شدیم}.
\end{quote}

\begin{quote}
\emph{Az khiyâbân-e asli vâred-e yek khiyâbân-e far'i shodim.}
\end{quote}

\begin{quote}
From the main road we turned into a side street.
\end{quote}

The opposite of \emph{asli}.

\begin{cousinweb}
\emph{far'i} was borrowed from Arabic and is usually traced to Arabic \emph{farʿī} 'of a branch', from \emph{farʿ} 'branch'. EtymDB does not list it.
\end{cousinweb}

\begin{grammarlens}[title={What you've built}]
One more word for telling, explaining and describing at B1.
\end{grammarlens}

\subsection*{Your turn}
\begin{itemize}
\raggedright
  \item \emph{Say it:} \emph{far'i}
  \item \emph{Say it:} \emph{far'i}, once more
  \item \emph{Say it:} a sentence of your own with \emph{far'i}
  \item \emph{Recall:} say the Persian for a postcode, then the Persian for main, then say \emph{far'i} again
\end{itemize}

\begin{tcolorbox}[breakable,colback=teal!4,colframe=teal!35!black,title={Before you move on}]
What does \fa{فرعی} mean? (\textquotedblleft{}Side, minor (of a road); secondary\textquotedblright{}.) Say it once more.
\end{tcolorbox}
